% Generated draft: requires mathematical and source review.
\documentclass[11pt]{article}
\usepackage[margin=25mm]{geometry}
\usepackage{fontspec}
\setmainfont{Latin Modern Roman}
\usepackage{amsmath,amssymb,mathtools,unicode-math}
\setmathfont{Latin Modern Math}
\usepackage{enumitem,needspace,xurl,hyperref}
\setlength{\parindent}{0pt}
\setlength{\parskip}{5pt}
\setlength{\emergencystretch}{3em}
\newcommand{\N}{\mathbb{N}}
\newcommand{\Z}{\mathbb{Z}}
\newcommand{\Q}{\mathbb{Q}}
\newcommand{\R}{\mathbb{R}}
\newcommand{\C}{\mathbb{C}}
\newcommand{\F}{\mathbb{F}}
\newcommand{\A}{\mathbb{A}}
\newcommand{\PP}{\mathbb{P}}
\newcommand{\E}{\mathbb{E}}
\newcommand{\eps}{\varepsilon}
\newcommand{\dd}{\,\mathrm{d}}
\begin{document}
\section*{UnsolvedMath ID: 30004435}
\textbf{Source problem code:} catalogue entry
\subsection{Definitions and mathematical statement}
Let \(S\) be a finite nonempty set equipped with its discrete sigma-field, and let \((X_k)_{k\in\Z}\) be an \(S\)-valued process on a probability space \((\Omega,\mathcal F,\mathbb P)\). Assume strict stationarity: for every \(r\geq1\), every \(k_1,\ldots,k_r\in\Z\), and every \(h\in\Z\),
\[
(X_{k_1+h},\ldots,X_{k_r+h})\overset{\mathrm d}= (X_{k_1},\ldots,X_{k_r}).
\]
Define the left and right tail sigma-fields by
\[
\mathcal T^-:=\bigcap_{n=0}^{\infty}\sigma(X_k:k\leq -n),
\qquad
\mathcal T^+:=\bigcap_{n=0}^{\infty}\sigma(X_k:k\geq n).
\]
Equality modulo null sets means that for every \(A\in\mathcal T^-\) there exists \(B\in\mathcal T^+\) with \(\mathbb P(A\triangle B)=0\), and conversely for every \(B\in\mathcal T^+\) there exists such an \(A\in\mathcal T^-\).

\textbf{Requested statement.} For every process satisfying these assumptions, prove, using only probability theory and without entropy theory, that
\[
\mathcal T^-=\mathcal T^+\qquad\text{modulo }\mathbb P\text{-null sets}.
\]
\textit{Unresolved qualifications.} The supplied statement does not explicitly specify the time index set or the meaning of stationarity. The formulation above uses the conventional interpretation of a bi-infinite discrete-time strictly stationary process; this interpretation requires source confirmation. The source also does not formally delimit the permitted probabilistic methods or the excluded entropy-theoretic methods. Its target is a proof subject to this methodological restriction.
\subsection{Short English statement}
For a stationary process taking values in a finite set, show that the events determined by the arbitrarily remote past coincide, up to probability-zero differences, with those determined by the arbitrarily remote future. The requested proof must use only probability theory and exclude entropy theory. Bi-infinite discrete time and strict stationarity are the explicit interpretations used here; the source leaves the precise methodological boundary unspecified.
\subsection{Sources}
\begin{itemize}
\item Supplied local registry record: UnsolvedMath permanent ID \(30004435\), title \emph{Probabilistic Equality of Left and Right Tail Fields}, exact source problem code \texttt{catalogue entry}, in \emph{Oberwolfach Reports Open Problems}. No individual proposer is identified in the supplied record.
\item Supplied report citation: \emph{Mini-Workshop: One-sided and Two-sided Stochastic Descriptions} (2021), Oberwolfach Reports, DOI 10.4171/owr/2020/11, pp.\ 601--637. Supplied source URL: \url{https://doi.org/10.4171/owr/2020/11}. The record gives proposed year 2020 and an \texttt{owr-year:2020} tag; their relationship to the citation's 2021 date is unresolved here.
\item Supplied literature source, labelled \texttt{OWR-2020-11}: \url{https://ems.press/content/serial-article-files/46847}.
\item Literature assessment dated exactly \texttt{2026-08-22}: the supplied curation states that equality of the two tails is known, but that no proof using only probability theory was identified, so it labels the method-restricted problem open. This is supplied dated curation, not an independently verified literature assessment; the background expressly requests independent verification of time-sensitive status and cited claims.
\item The registry supplies status \texttt{open} and statement status \texttt{corrected\_verified}, with a statement-verification note reporting manual semantic, typographic, and LaTeX review. These are supplied registry labels and review claims, not independent mathematical verification of this draft.
\end{itemize}
\end{document}
